% !TeX root = ../main.tex
\section{Introduction}
\label{sec:introduction}

Probabilistic programming provides a flexible way to express computations involving random choices using familiar programming constructs~\cite{kozen1981semantics,gordon2014probabilistic,vandeMeent2018introduction,carpenter2017stan,bingham2019pyro,cusumano2019gen}. The combination of random sampling with features such as conditionals, higher-order functions, and recursion makes probabilistic programs expressive, but also challenging to analyze. A fundamental quantitative property of interest is the expected value of a program's output.

Existing approaches to computing expected values broadly follow two directions. Exact inference and moment analysis compute expectations symbolically for supported classes of programs~\cite{gehr2016psi,gehr2020lpsi,moosbrugger2022moment}. Alternatively, Monte Carlo methods estimate expectations through repeated sampling. Conditional Monte Carlo and Rao-Blackwellization improve the precision of these estimates by analytically averaging over selected random choices~\cite{casella1996rao,owen2013montecarlo}. Other techniques combine symbolic reasoning with sampling during inference~\cite{murray2018delayed,atkinson2022semisymbolic}, while expectation programming directly targets expected value estimation~\cite{reichelt2022expectation}. Together, these approaches illustrate that random choices need not always be explicitly sampled to compute or estimate expectations, suggesting the possibility of simplifying probabilistic programs themselves by eliminating selected random choices while preserving their expected values. Identifying which choices can be safely eliminated, however, is challenging: sampled values may influence control flow, determine the parameters of subsequent distributions, or participate in nonlinear computations. This raises a natural question:
\begin{center}
\emph{Can static analysis automatically identify and eliminate
random choices from a probabilistic program while preserving
its expected value?}
\end{center}

Consider the program in
\cref{fig:intro-signal}, which samples a signal $x$ uniformly
from $[0,1]$, then samples a Gaussian measurement $y$ with mean $x$ and a Gaussian measurement $z$ with mean $y$,
then returns the product $xz$. Once $x$ is fixed, the measurement
noise contributes nothing to the average product:
$$
  \mathbb E[xz\mid x]
  =x\,\mathbb E[z\mid x]
  =x\,\mathbb E[\mathbb E[z\mid x,y]\mid x]
  =x\,\mathbb E[y\mid x]
  =x^2.
$$
Replacing both Gaussian samples by their means therefore yields a simpler
program that samples only $x$ and returns $x^2$. Both programs have
expected value $1/3$, but the transformed program avoids generating
measurement noise and thus has lower output variance (\Cref{fig:intro-variance}).

% !TeX root = ../main.tex
\begin{figure}[htbp]
\centering
\begin{minipage}[t]{.48\linewidth}
\raggedright\textbf{Source (annotations inferred)}
\begin{lstlisting}[xleftmargin=0pt]
let x = uniform[G](0, 1) in
let y = gaussian[E](x, 1) in
let z = gaussian[E](y, 1) in
x * z
\end{lstlisting}
\end{minipage}\hfill
\begin{minipage}[t]{.48\linewidth}
\raggedright\textbf{After replacement (simplified)}
\begin{lstlisting}[xleftmargin=0pt]
let x = uniform[G](0, 1) in
x * x
\end{lstlisting}
\end{minipage}
\caption{A signal measured through two stages of noise.}
\label{fig:intro-signal}
\Description{The first sample remains random and both Gaussian samples are replaced
by their means. Both y and z become x in the simplified program.
The output becomes x squared, preserving expectation one third.}
\end{figure}


This replacement also reveals the central difficulty. Replacing the
remaining uniform sample by its mean, $1/2$, would make the program
return $1/4$ rather than $1/3$. Thus, only two random choices can be removed here, and
the third must remain. The distinction depends on how their values
are used: the returned product is linear in $z$ when $x$ is fixed,
but becomes quadratic in $x$ after averaging over $y$ and $z$. A sound
transformation must therefore reason about both the dependencies
between random choices and their uses throughout the program. This challenge is particularly pronounced in higher-order, potentially unbounded recursive programs, where the influence of a random choice can propagate through function calls and across arbitrarily long executions. 

Our key insight is that safe replacements can nevertheless be
identified compositionally from how sampled values are used,
without requiring symbolic marginalization of their distributions.
We introduce \emph{determinization}, a type-directed,
source-to-source transformation of higher-order, recursive
probabilistic programs with continuous distributions that
automatically identifies which samples can be replaced by
their means. Our type system distinguishes sampling choices
that remain random ($\G$) from those that can be eliminated
($\E$). Type inference determines these modes by tracking
how sampled values are used: values marked $\E$ may
participate in affine computations, but cannot determine
control flow or the parameters of retained samples.
The transformation then replaces each $\E$-mode sample by
its mean. We establish soundness through a
relational proof that connects the original and transformed
programs via two interpretations of a common symbolic
semantics. Under suitable conditions, the resulting program
preserves the expected value of the output while never
increasing its variance.

The next three examples illustrate, at a high level, how determinization handles
higher-order functions, data structures, and unbounded recursion.
We begin with the Gaussian random walk in
\cref{fig:intro-gaussian-walk}, from~\cite{vakar2019domain}.
The program passes a probabilistic step function to
\texttt{reduce}, a left fold that traverses the requested times
and accumulates a list of time-position pairs.
At each time $t'$, it samples a new position from
$\mathcal N(y,t'-t)$, where $(t,y)$ is the latest pair in the path. It then adds $(t',y')$ to the front of the path. At the next step, it reads that pair and uses $y'$ as the next Gaussian's mean. Inference marks the Gaussian sampling
site as determinizable without unrolling the traversal.
Replacing the draws by their means keeps every position at zero,
yielding a deterministic path that matches the expected position
at each requested time. 

% !TeX root = ../main.tex
\begin{figure}[htbp]
\centering
\begin{minipage}[t]{.48\linewidth}
\raggedright\textbf{Source (annotations inferred)}
\begin{lstlisting}[xleftmargin=0pt]
let step = fun path => fun t' =>
  match path with
    | [] => []
    | current :: _ =>
        let t = fst current in
        let y = snd current in
        let y' = gaussian[E](y, t' - t) in
        (t', y') :: path
in
let times =
  1 :: 2 :: 4 :: 5.5 :: 6 :: 6.25 :: []
in
reduce step ((0, 0) :: []) times
\end{lstlisting}
\end{minipage}\hfill
\begin{minipage}[t]{.48\linewidth}
\raggedright\textbf{After replacement (simplified)}
\begin{lstlisting}[xleftmargin=0pt]
let step = fun path => fun t' =>
  match path with
    | [] => []
    | current :: _ =>
        let y = snd current in
        (t', y) :: path
in
let times =
  1 :: 2 :: 4 :: 5.5 :: 6 :: 6.25 :: []
in
reduce step ((0, 0) :: []) times
\end{lstlisting}
\end{minipage}
\caption{Gaussian random walk from~\cite{vakar2019domain}, consisting of a higher order function \texttt{reduce} and a recursive traversal of lists of pairs. Full program in~\cref{sec:appendix-gauss-random-walk}.}
\label{fig:intro-gaussian-walk}
\Description{The source passes a probabilistic step function to the recursive
reduce function. At each requested time, the step draws a Gaussian position
whose mean is the previous position and whose variance is the elapsed time, and
prepends the new time--position pair to the path. Replacing every draw by its
mean leaves the position fixed at zero.}
\end{figure}


Beyond the transformation itself, determinization can improve
Monte Carlo estimation of expected values.
Consider the dungeon crawl program in \cref{fig:intro-dungeon}, which recursively visits
rooms until an exit coin ends the computation, with no fixed
bound on the number of rooms visited.
At depth $d$, a room pays a jackpot of $200d$ with probability
$0.01$. The exit coins remain random because they control
recursion, while each jackpot contributes linearly to the total
reward and is replaced by its mean payout, $2d$. The determinized program preserves the expected value while reducing output
variance by a factor of $12.36$ (\cref{fig:intro-variance}), hence requiring correspondingly fewer samples to achieve the same estimator variance.

% !TeX root = ../main.tex
\begin{figure}[htbp]
\centering
\begin{minipage}{\linewidth}
\raggedright
\lstset{xleftmargin=0pt}
\textbf{Dungeon crawl: source (annotations inferred)}
\begin{lstlisting}
let crawl = fun loot =>
  (rec visit depth =>
    if flip[G](0.25) then 0
    else loot depth + visit (depth + 1)
  ) 1
in
crawl (fun depth => 200 * depth * bernoulli[E](0.01))
\end{lstlisting}
\medskip
\textbf{After replacement (simplified)}\quad\texttt{crawl} unchanged
\begin{lstlisting}
crawl (fun depth => 2 * depth)
\end{lstlisting}
\end{minipage}
\caption{Dungeon crawl: \texttt{crawl} is a higher order function while \texttt{visit} uses unbounded recursion.}
\label{fig:intro-dungeon}
\Description{The complete source program defines a recursive crawl and passes
it a random loot function. Exit coins remain random; jackpot coins become their
means. Only the final call changes, to a loot function returning twice the depth.}
\end{figure}


% !TeX root = ../main.tex
\begin{figure}[htbp]
\centering
\begin{tikzpicture}[font=\footnotesize]
\begin{groupplot}[
  group style={group size=2 by 1, horizontal sep=1.6cm},
  width=.43\linewidth, height=4.6cm,
  axis line style={black!45}, tick style={black!45},
  grid=major, grid style={black!10},
  title style={font=\small\bfseries},
  ylabel style={font=\footnotesize}, xlabel style={font=\footnotesize},
]
\nextgroupplot[title={Signal: output variance},
  ymin=0,ymax=.6,ytick={0,.2,.4,.6},
  xmin=-.6,xmax=1.6,xtick={0,1},xticklabels={Source,Determinized},
  scaled y ticks=false]
\addplot[ybar,bar width=25pt,draw=black!65,fill=black!18]
  coordinates {(0,\signalSourceVarianceValue)};
\addplot[ybar,bar width=25pt,draw=teal!75!black,fill=teal!25]
  coordinates {(1,\signalTargetVarianceValue)};
\node[anchor=south] at (axis cs:0,\signalSourceVarianceValue) {$\displaystyle\signalSourceVariance$};
\node[anchor=south] at (axis cs:1,\signalTargetVarianceValue) {$\displaystyle\signalTargetVariance$};
\node[anchor=north east,text=teal!75!black] at (rel axis cs:.98,.98)
  {$\signalReduction\times$ reduction};

\nextgroupplot[title={Dungeon: output variance},
  ymin=0,ymax=48000,ytick={0,10000,20000,30000,40000},
  xmin=-.6,xmax=1.6,xtick={0,1},xticklabels={Source,Determinized},
  scaled y ticks=false]
\addplot[ybar,bar width=25pt,draw=black!65,fill=black!18]
  coordinates {(0,\dungeonSourceVarianceValue)};
\addplot[ybar,bar width=25pt,draw=teal!75!black,fill=teal!25]
  coordinates {(1,\dungeonTargetVarianceValue)};
\node[anchor=south] at (axis cs:0,\dungeonSourceVarianceValue) {$\dungeonSourceVariance$};
\node[anchor=south] at (axis cs:1,\dungeonTargetVarianceValue) {$\dungeonTargetVariance$};
\node[anchor=north east,text=teal!75!black] at (rel axis cs:.98,.98)
  {$\dungeonReduction\times$ reduction};
\end{groupplot}
\end{tikzpicture}
\caption{Analytically calculated output variances for the signal example
(\cref{fig:intro-signal}) and dungeon crawl (\cref{fig:intro-dungeon}).
Averaging $n$ independent runs divides variance by $n$,
so the reduction factors also apply to the runs needed for the same variance of the average.}
\label{fig:intro-variance}
\Description{Source and determinized output variances are 19/45 and 4/45 for
the signal, and 36192 and 2928 for the dungeon. The variance reduction factors,
4.75 and approximately 12.36, also give the reduction in runs needed for the
same variance of the sample average.}
\end{figure}


Moreover, removing randomness can do more than improve sampling: it can enable exact expectation computation by transforming an infinite-state program into a finite-state one. In \cref{fig:intro-exact}, a recursive process either returns
its current state $x$ or continues from $1-x$ plus Gaussian
noise, each with probability $1/2$. Determinization removes the noise: starting from zero, the transformed loop has only two possible states, $0$ and $1$.
Let $v_i$ denote the expected value eventually returned by the
transformed program when entering the loop with $\texttt{x}=i$,
accounting for all subsequent iterations. The expected outputs
from these two states satisfy
\[
 \begin{pmatrix}v_0\\v_1\end{pmatrix}
 =\begin{pmatrix}0&\tfrac12\\\tfrac12&0\end{pmatrix}
  \begin{pmatrix}v_0\\v_1\end{pmatrix}
  +\begin{pmatrix}0\\\tfrac12\end{pmatrix}
 \qquad\Longrightarrow\qquad
 \begin{pmatrix}v_0\\v_1\end{pmatrix}
 =\begin{pmatrix}\tfrac13\\\tfrac23\end{pmatrix}.
\]
Since the program starts with $\texttt{x}=0$, its expected
output is $v_0=\tfrac13$.
A continuous-state computation thus becomes a finite-state
problem whose expectation can be computed exactly. Our implementation produces a Lean-checkable certificate for the resulting exact value.

% !TeX root = ../main.tex
\begin{figure}[htbp]
\centering
\begin{minipage}[t]{.48\linewidth}
\raggedright\textbf{Source (annotations inferred)}
\begin{lstlisting}[xleftmargin=0pt]
(rec run x =>
  if flip[G](0.5) then x
  else run
    (1 - x + gaussian[E](0, 1))
) 0
\end{lstlisting}
\end{minipage}\hfill
\begin{minipage}[t]{.48\linewidth}
\raggedright\textbf{After replacement (simplified)}
\begin{lstlisting}[xleftmargin=0pt]
(rec run x =>
  if flip[G](0.5) then x
  else run (1 - x)
) 0
\end{lstlisting}
\end{minipage}
\par\medskip
Let $v_i$ be the expected output from loop entry $\texttt{x}=i$
in the transformed program.
\[
 \begin{pmatrix}v_0\\v_1\end{pmatrix}
 =\begin{pmatrix}0&\tfrac12\\\tfrac12&0\end{pmatrix}
  \begin{pmatrix}v_0\\v_1\end{pmatrix}
  +\begin{pmatrix}0\\\tfrac12\end{pmatrix}
 \qquad\Longrightarrow\qquad
 \begin{pmatrix}v_0\\v_1\end{pmatrix}
 =\begin{pmatrix}\tfrac13\\\tfrac23\end{pmatrix}.
\]
\caption{A noisy iteration becomes a finite-state loop.
Both programs have expected output $v_0=1/3$.}
\label{fig:intro-exact}
\Description{The source repeatedly replaces x by one minus x plus Gaussian
noise, stopping with probability one half at each iteration. Removing the noise
leaves two possible loop-entry values, zero and one. Their expected outputs
satisfy v0 equals v1 over two, and v1 equals one half plus v0 over two,
giving initial expected output one third.}
\end{figure}


\paragraph{Contributions.}
\begin{itemize}
  \item \textbf{Type system and inference.}
  We develop a type system and source-to-source
  transformation that identifies and replaces eligible samples by their means, even in higher-order,
  potentially unbounded recursive programs.
  Our inference algorithm finds the greatest assignment of
  determinizable sites permitted by the type system,
  without requiring symbolic marginalization or unrolling
  recursive computations
  (\cref{sec:language,sec:type-system}).

  \item \textbf{Soundness.} We prove that determinization preserves expected values
  and never increases output variance, under suitable
  assumptions. Since the transformation changes the output
  distribution, its correctness requires relating the
  original and transformed executions rather than showing
  distributional equivalence. We develop a relational proof
  based on two interpretations of a common symbolic
  semantics, establishing a stronger tracewise soundness
  result from which the global guarantees follow
  (\cref{sec:soundness,sec:symbolic-semantics}).

    \item \textbf{Lean implementation and verification.}
    We implement type inference, determinization, and sampling
  in Lean, and mechanize our soundness theorems.
  Our implementation additionally supports verified
  matrix solving and certificate checking, providing an
  end-to-end correctness guarantee: if the original program
  is domain-safe and its expectation is well-defined, the
  computed expectation is exactly that of the original
  program
  (\cref{sec:implementation}).
  The full Lean code is available at [TODO], with references
  to mechanized definitions and theorems indicated by a
  superscript Lean logo, as in \leanmark.

    \item \textbf{Evaluation.} We evaluate two applications of determinization.
First, by eliminating continuous sampling, determinization
can reduce infinite-state programs to finite-state programs,
enabling exact computation of their expected values using
existing model-checking techniques.
Second, even when the transformed program remains infinite-state,
reducing its randomness can lower output variance and improve
the precision of Monte Carlo estimates under a fixed
wall-clock budget (\cref{sec:evaluation}).
\end{itemize}
